% Generated draft: requires mathematical and source review.
\documentclass[11pt]{article}
\usepackage[margin=25mm]{geometry}
\usepackage{fontspec}
\setmainfont{Latin Modern Roman}
\usepackage{amsmath,amssymb,mathtools,unicode-math}
\setmathfont{Latin Modern Math}
\usepackage{enumitem,needspace,xurl,hyperref}
\setlength{\parindent}{0pt}
\setlength{\parskip}{5pt}
\setlength{\emergencystretch}{3em}
\newcommand{\N}{\mathbb{N}}
\newcommand{\Z}{\mathbb{Z}}
\newcommand{\Q}{\mathbb{Q}}
\newcommand{\R}{\mathbb{R}}
\newcommand{\C}{\mathbb{C}}
\newcommand{\F}{\mathbb{F}}
\newcommand{\A}{\mathbb{A}}
\newcommand{\PP}{\mathbb{P}}
\newcommand{\E}{\mathbb{E}}
\newcommand{\eps}{\varepsilon}
\newcommand{\dd}{\,\mathrm{d}}
\begin{document}
\section*{UnsolvedMath ID: 30000591}
\textbf{Source problem code:} catalogue entry
\subsection{Definitions and mathematical statement}
Over $\F_2$, let $(V,q)$ be a quadratic space and let $O(V,q)$ be the group of linear automorphisms preserving $q$. The polar form is $b_q(v,w)=q(v+w)-q(v)-q(w)$. Preserving $q$, not merely its polar form, specifies the orthogonal group in characteristic two. The source does not provide the dimensions, nondegeneracy conditions, or plus/minus type of the spaces in its stabilization system.

The first question asks whether the intended inverse system of orthogonal-group cohomology stabilizes. To express the usual degreewise meaning, suppose the report gives spaces $(V_r,q_r)$, compatible group inclusions, coefficient modules $M_r$, and cohomological restriction maps
\[
 \cdots\longrightarrow H^i(O(V_{r+1},q_{r+1});M_{r+1})
       \longrightarrow H^i(O(V_r,q_r);M_r)\longrightarrow\cdots.
\]
Stabilization would mean that for each relevant degree $i$ the transition maps are isomorphisms for all sufficiently large $r$, with any coefficient data fixed in the report's sense. The threshold may depend on $i$ and the coefficient functor. The extracted statement gives none of these modules, maps, or quantifiers, so this is an explanation of the missing system rather than an assertion for ordinary trivial-coefficient cohomology.

The second question asks whether the stable value is described by an analogue of the Betley--Suslin theorem in $\mathcal F_{\mathrm{quad}}$. That category is named but not defined. Its source category of quadratic objects and morphisms, target category, allowed functors, and comparison construction must be recovered. In particular, it must not be replaced without evidence by ordinary functors on vector spaces or merely on isometric embeddings. The source also does not supply the version of the Betley--Suslin theorem, the intended Ext/cohomology comparison, or the proposed formula for the stable value.

The questions concern both eventual stabilization and identification of the resulting stable object. Existence of an inverse limit alone is not the requested eventual-isomorphism property, and neither part is silently replaced by the analogous general-linear-group result.
\subsection{Short English statement}
The entry asks whether a specified orthogonal-group cohomology system over $\F_2$ eventually stabilizes and whether its stable value admits a functor-category description. The coefficient system, stabilization maps, and quadratic-functor comparison are missing.
\subsection{Sources}
\begin{itemize}
\item Local UnsolvedMath registry, permanent ID 30000591, code catalogue entry, ``Stable Cohomology of Orthogonal Groups over Binary Fields.'' Proposed year 2006; proposer unspecified.
\item \emph{Topologie} (2007), Oberwolfach Reports, pp.\ 2579--2630, DOI 10.4171/owr/2006/43; \url{https://doi.org/10.4171/owr/2006/43}. The displayed report year 2007 differs from the DOI and proposed year 2006.
\item Supplied reference ``The functor category Fquad'': \url{https://arxiv.org/abs/math/0606484}.
\item Statement status corrected\_verified; the supplied note says the field and functor-category notation was enclosed in balanced delimiters, preserving both questions.
\item Literature triage checked 2026-08-21: open. It says the stabilization and Betley--Suslin-analogue questions remain research questions and the cited papers develop the category without giving the requested general answers. Exact hypotheses and current status require independent verification. Created 2026-08-25; updated 2026-09-16.
\item Registry curation by Ulam AI / UnsolvedMath Contributors, CC BY 4.0; underlying report and literature provenance remain as supplied. Unreviewed statement-only draft.
\end{itemize}
\end{document}
